% Generated draft: requires mathematical and source review.
\documentclass[11pt]{article}
\usepackage[margin=25mm]{geometry}
\usepackage{fontspec}
\setmainfont{Latin Modern Roman}
\usepackage{amsmath,amssymb,mathtools,unicode-math}
\setmathfont{Latin Modern Math}
\usepackage{enumitem,needspace,xurl,hyperref}
\setlength{\parindent}{0pt}
\setlength{\parskip}{5pt}
\setlength{\emergencystretch}{3em}
\newcommand{\N}{\mathbb{N}}
\newcommand{\Z}{\mathbb{Z}}
\newcommand{\Q}{\mathbb{Q}}
\newcommand{\R}{\mathbb{R}}
\newcommand{\C}{\mathbb{C}}
\newcommand{\F}{\mathbb{F}}
\newcommand{\A}{\mathbb{A}}
\newcommand{\PP}{\mathbb{P}}
\newcommand{\E}{\mathbb{E}}
\newcommand{\eps}{\varepsilon}
\newcommand{\dd}{\,\mathrm{d}}
\begin{document}
\section*{UnsolvedMath ID: 20000625}
\textbf{Source problem code:} catalogue entry
\subsection{Definitions and mathematical statement}
This source record concatenates a collection of spectral questions and historical remarks, rather than one isolated problem. The formulas below distinguish its reconstructed notation, source-time literature reports, and the separate correction-transfer proposition.

Let \(\Omega\subset\R^n\), \(n\geq2\), be a nonempty bounded open set. Let \(0<E_1\leq E_2\leq\cdots\) be the Dirichlet eigenvalues of \(-\Delta\), defined through \(H^1_0(\Omega)\), with multiplicity. Put \(\omega_n=|S^{n-1}|/n\),
\[
 W_\Omega=(2\pi)^2(\omega_n|\Omega|)^{-2/n},\quad
 R_\sigma(\Lambda)=\sum_j(\Lambda-E_j)_+^\sigma,\quad
 L_{\sigma,n}^{\mathrm{cl}}=\frac{\Gamma(\sigma+1)}{(4\pi)^{n/2}\Gamma(\sigma+1+n/2)}.
\]
For \(\sigma=0\), use the strict counting convention \((\Lambda-E_j)_+^0=1_{E_j<\Lambda}\).

The first section contains the following questions and background.
\begin{enumerate}
\item Michael Loss and Timo Weidl: P\'olya's conjecture is \(E_j\geq W_\Omega j^{2/n}\). The source says it remains open even at \(j=3\); this is historical wording. It recalls Li--Yau,
\[
 \sum_{j=1}^JE_j\geq\frac n{n+2}W_\Omega J^{1+2/n},
\]
and Berezin's 1972 bound \(R_\sigma(\Lambda)\leq L_{\sigma,n}^{\mathrm{cl}}|\Omega|\Lambda^{\sigma+n/2}\) for \(\sigma\geq1\). It asks for \(0\leq\sigma<1\), with \(\sigma=0\) equivalent to P\'olya, and recalls that order one yields Li--Yau by Legendre duality.
\item Timo Weidl: seek lower-order corrections to Li--Yau, especially involving boundary area. The source reports discrete-domain corrections, Melas's moment-of-inertia correction, a February 2009 two-dimensional boundary correction by Kovarik--Vugalter--Weidl, and Weidl's earlier effective-boundary improvement for Riesz order at least \(3/2\). It describes further improvements below \(3/2\) as desirable; no universal boundary-profile formula or complete domain hypotheses are supplied.
\item Timo Weidl: seek corrections to the reversed Neumann Li--Yau analogue of Kr\"oger. The source describes Neumann P\'olya as open for \(j\geq2\), apart from a then-recent two-dimensional \(j=2\) result of Girouard and others. Neumann indexing, regularity and compact-resolvent assumptions are not specified in the excerpt, so no unverified indexing formula is imposed.
\item Timo Weidl: analogous questions in magnetic fields, referring to an unavailable later item (10).
\item Evans Harrell and Joachim Stubbe: for \(\sigma\geq2\), the source reports that
\[
 r_\sigma(\Lambda)=\Lambda^{-\sigma-n/2}
 \bigl(L_{\sigma,n}^{\mathrm{cl}}|\Omega|\Lambda^{\sigma+n/2}-R_\sigma(\Lambda)\bigr)
\]
is nonincreasing and tends to zero at infinity. It asks for boundary corrections from the trace-identity method. For periodic boundary conditions it reports a related monotonicity and a connection with the Gauss circle or lattice-point problem. The extracted dimension \(d\) is reconstructed as \(n\), and the phase-space integral uses a positive part.
\item Harrell and Stubbe: seek analogous monotonicity and Berezin--Li--Yau results for higher-order operators, such as the clamped plate, and fractional Laplacians; a then-accepted Klein--Gordon paper is mentioned for \(\sqrt{-\Delta}\). Precise realizations and boundary conditions beyond the example are not supplied.
\item Harrell and Stubbe: define
\[
 M_p(J)=\left(\frac{n+2p}{nJ}\sum_{j=1}^JE_j^p\right)^{1/p}\ (p>0),\qquad
 M_0(J)=e^{2/n}\left(\prod_{j=1}^JE_j\right)^{1/J}.
\]
All have Weyl leading term \(W_\Omega J^{2/n}\). The supplied verification explicitly reconstructs the damaged dispersion as \(\Delta_J=M_1(J)^2-M_2(J)^2\), not the dimensionally inconsistent printed \(M_1^2-M_2\). The reported inequalities are
\[
 \Delta_J\geq\tfrac14(E_{J+1}-E_J)^2,\qquad
 M_1(J)-\sqrt{\Delta_J}\leq E_J\leq E_{J+1}\leq M_1(J)+\sqrt{\Delta_J}.
\]
For \(p>0\), seek a bound
\[
 M_p(J)^{2p}-M_{2p}(J)^{2p}
 \leq C(p,\Omega)E_1^{2p}J^{2p\kappa},\qquad\kappa<2/n,
\]
with the left side also an explicit reconstruction. The permitted dependence of \(\kappa\) is not stated.
\item Harrell and Stubbe: ask whether \(E_J\leq M_p(J)\) for every \(J\geq1\), \(p\geq0\), and whether some \(\Omega,J\) attain equality in the gap-dispersion inequality.
\end{enumerate}

The supplied aid addresses a correction-transfer part of these questions. Let \(q=1+n/2\), \(A=L_{1,n}^{\mathrm{cl}}|\Omega|\), and
\[
 \Lambda_J=\left(\frac{J}{qA}\right)^{1/(q-1)}=W_\Omega J^{2/n}.
\]
For any domain-dependent profile \(\Phi_\Omega\), if the single-cutoff bound
\[
 R_1(\Lambda_J)\leq A\Lambda_J^q-\Phi_\Omega(\Lambda_J)
\]
holds, then the aid asserts the exact implication
\[
 M_1(J)\geq\Lambda_J+\frac{n+2}{nJ}\Phi_\Omega(\Lambda_J).
\]
No global-in-cutoff estimate is required for this implication at a fixed \(J\). In contrast, propagation to higher orders uses an order-one bound throughout the integration interval: for \(\sigma>1\),
\[
 R_\sigma(\Lambda)=\sigma(\sigma-1)\int_0^\Lambda(\Lambda-t)^{\sigma-2}R_1(t)\,\dd t,
\]
and the correction becomes \(\Phi_{\Omega,\sigma}(\Lambda)=\sigma(\sigma-1)\int_0^\Lambda(\Lambda-t)^{\sigma-2}\Phi_\Omega(t)\,\dd t\). Where a valid order-one profile has boundary coefficient \(\tfrac14|\partial\Omega|L_{1,n-1}^{\mathrm{cl}}\), this integration carries it to \(\tfrac14|\partial\Omega|L_{\sigma,n-1}^{\mathrm{cl}}\). This is conditional on the profile's validity, not a universal boundary correction for every domain and cutoff. Applied to the reported proved two-term Dirichlet Riesz asymptotic for \(\Omega=(0,1)^2\), the aid obtains
\[
 M_1(J)=4\pi J+\frac{32\sqrt\pi}{3}J^{1/2}+o(J^{1/2})\qquad(J\to\infty).
\]

The record also includes a second section and extended remarks on Lieb--Thirring problems. For \(n\geq1\), \(V\geq0\), and the negative eigenvalues \(\lambda_j\) of \(-\Delta-V\) on \(L^2(\R^n)\), use
\[
 \sum_j|\lambda_j|^\gamma\leq L_{n,\gamma}\int_{\R^n}V^{\gamma+n/2},\qquad
 C_{n,\gamma}=L_{n,\gamma}/L_{\gamma,n}^{\mathrm{cl}}.
\]
The source lists the ranges \(\gamma\geq1/2\) for \(n=1\), \(\gamma>0\) for \(n=2\), and \(\gamma\geq0\) for \(n\geq3\), with the last zero-order case being CLR. It reports \(C_{n,\gamma}\geq1\) and Aizenman--Lieb monotonicity in \(\gamma\). Its historical table records: \(C_{1,1/2}=2\); for \(1/2<\gamma<3/2\), upper bound 2 with improvements for some exponents and conjectured optimum
\(2((\gamma-1/2)/(\gamma+1/2))^{\gamma-1/2}\); and \(C_{n,\gamma}=1\) for \(\gamma\geq3/2\). The February 2009 entries give 3.64 for \(1/2\leq\gamma<1\), \(n\geq2\), and 1.82 for \(1\leq\gamma<3/2\), \(n\geq2\), with optimum 1 conjectured in the latter range for \(n\geq3\). For \(0\leq\gamma<1/2\), the table gives 6.87 for \(n=3\), conjectured \(8/\sqrt3\), and 10.34 for \(n\geq4\); for \(n=2\), \(0<\gamma<1/2\), it prints a question mark. Its Eden--Foias note reads numerically about 1.82 at \((n,\gamma)=(1,1)\), consistent with \(\pi/\sqrt3\); the extracted fraction is damaged. These are source-time bounds and conjectures, not verified present-day best constants.

The appended questions are: Richard Laugesen asks whether an optimal potential exists, particularly for \(n=1\), \(1/2<\gamma<3/2\), including the fixed-number-of-bound-states variant, and whether an optimizer can have a single bound state. The source relates this to Keller's one-eigenvalue optimizer and the 1976 Lieb--Thirring conjecture, and mentions endpoint delta and reflectionless potentials; admissible function versus measure classes need clarification. Eric Carlen asks for \(\sum_j|\lambda_j|^\gamma\leq C|\lambda_1|^\gamma\), allowing dependence on dimension, exponent, and unspecified integrability data. Rafael Benguria asks for Lieb--Thirring analogues for the Benjamin--Ono Lax operator, noting workshop progress by Ekholm, Frank and Hundertmark on monotonicity. Rupert Frank asks for a direct proof and corrections to the Eden--Foias bound; the February 2009 note reports an operator-valued extension and dimension lifting. Timo Weidl asks for direct eigenvalue-sum estimates without passing through Birman--Schwinger counting.

Almut Burchard's appended Ovals Problem concerns a smooth closed curve \(\gamma\subset\R^3\) of length \(2\pi\), curvature \(\kappa(s)\), and
\[
 H_\gamma=-\frac{\dd^2}{\dd s^2}+\kappa(s)^2
\]
on \(2\pi\)-periodic functions. It asks whether the principal eigenvalue \(e(\gamma)\geq1\). The source says this is open even for convex planar loops, reports equality for a family with \(\kappa(s)=(a^2\cos^2s+a^{-2}\sin^2s)^{-1}\), \(a>0\), degenerating to two length-\(\pi\) segments as \(a\to0\), and records a then-best lower bound 0.6085, nearby-family results, and results under extra geometric constraints. It records discussions by Benguria, Loss, Burchard, Thomas and Linde; claims by Linde and Burchard about possibly cornered minimizers; a possible affine symmetry; proposed flows by Carlen, Mazzeo and Benguria, including affine curvature flow; ODE methods proposed by Rapti and Lee; and Laugesen's suggested Birman--Schwinger operator
\(T=\kappa(\dd^2/\dd s^2+c)^{-1}\kappa\), \(0<c<1\), whose largest eigenvalue is proposed to exceed one. The suggested endpoint \(c=1\) must address the kernel spanned by \(\sin s,\cos s\) and the stated restriction \(\kappa\psi\perp\sin s,\cos s\); strict inequality at equality curves is unresolved in the excerpt. The cornered-minimizer claims also require a relaxed admissible class. These ancillary questions and proposed approaches are preserved as source reports, not supplied solutions.
\subsection{Short English statement}
The record combines P\'olya, eigenvalue-sum, Lieb--Thirring, and Ovals questions with historical updates. The supplied partial theorem transfers an order-one Riesz deficit at one Weyl cutoff into a normalized eigenvalue-mean correction and propagates valid profiles to higher orders. It gives the explicit second term on the unit square; it does not solve the full collection.
\subsection{Sources}
\begin{itemize}
\item Local registry: permanent UnsolvedMath ID \(20000625\), exact public code \texttt{catalogue entry}; \texttt{aim-analysis-notes.json}, \texttt{notes[156]}. AIM workshop \emph{Low eigenvalues of Laplace and Schrodinger operators}, supplied item 1: \url{https://aimath.org/WWN/loweigenvalues/loweigenvalues.pdf}. This is supplied provenance; no PDF access is claimed.
\item J. K. Freericks, E. H. Lieb, D. Ueltschi, ``Phase separation due to quantum mechanical correlations,'' \emph{Phys. Rev. Lett.} 88, 106401, 1--4 (2002). A. Melas, ``A lower bound for sums of eigenvalues of the Laplacian,'' \emph{Proc. Amer. Math. Soc.} 131 (2003), 631--636.
\item H. Kovarik, S. Vugalter, T. Weidl, ``Two-dimensional Berezin-Li-Yau inequalities with a correction term,'' \emph{Comm. Math. Phys.} 287 (2009), 959--981. T. Weidl, ``Improved Berezin-Li-Yau inequalities with a remainder term,'' in \emph{Spectral theory of differential operators}, T. Suslina and D. Yafaev, eds., \emph{Amer. Math. Soc. Transl. Ser. 2}, 225 (2008).
\item P. Kr\"oger, ``Upper bounds for the Neumann eigenvalues on a bounded domain in Euclidean space,'' \emph{J. Funct. Anal.} 106 (1992), no. 2, 353--357. A. Laptev, ``Dirichlet and Neumann eigenvalue problems on domains in Euclidean spaces,'' \emph{J. Funct. Anal.} 151 (1997), 531--545. The source cites A. Girouard and others, \emph{J. Diff. Geometry}, ``to appear,'' without a title or complete citation.
\item E. M. Harrell and L. Hermi, ``Differential inequalities for Riesz means and Weyl-type bounds for eigenvalues,'' \emph{J. Funct. Analysis} 254 (2008), 3173--3191. E. M. Harrell and J. Stubbe, ``On trace identities and universal eigenvalue estimates for some partial differential operators,'' \emph{Trans. Amer. Math. Soc.} 349 (1997), 1797--1809.
\item E. M. Harrell and J. Stubbe, ``Universal bounds and semiclassical estimates for eigenvalues of abstract Schr\"odinger operators,'' preprint 2008, arXiv:0808.1133; and ``Trace identities for commutators, with applications to the distribution of eigenvalues,'' preprint 2009, arXiv:0903.0563. E. M. Harrell and S. Y\i ld\i r\i m Yolcu, ``Eigenvalue inequalities for Klein-Gordon Operators,'' described as accepted by \emph{J. Funct. Analysis}.
\item The appended remark cites lecture notes by Michael Loss and Timo Weidl and a survey by Dirk Hundertmark without complete details. February 2009 updates cite J. Dolbeault, A. Laptev, M. Loss, ``Lieb-Thirring inequalities with improved constants,'' \emph{J. Eur. Math. Soc.} 10 (2008), and R. L. Frank, E. H. Lieb, R. Seiringer, ``Number of bound states of Schr\"odinger operators with matrix-valued potentials,'' \emph{Lett. Math. Phys.} 82, 107 (2007).
\item J. B. Keller, \emph{J. Mathematical Phys.} 2, 262--266 (1961), title not supplied; the single-bound-state conjecture is attributed to Lieb and Thirring (1976). A. Eden and C. Foias, ``A simple proof of the generalized Lieb-Thirring inequalities of one-space dimension,'' \emph{Journal of mathematical analysis and applications} 162 (1991), 250--254.
\item R. L. Anderson and E. Tafflin, ``The Benjamin-Ono equation -Recursivity of linearization mapsLax pairs,'' \emph{Letters in Mathematical Physics} 9 (1985), 299--311, as supplied. D. J. Kaup and Y. Matsuno, ``The inverse scattering transform for the Benjamin-Ono equation,'' \emph{Studies in applied mathematics} 101 (1998), 73--98.
\item Ovals references [1]--[7] are invoked for curvature flows, narrow-channel quantum mechanics, Dirac operators, elastic-loop tension, local estimates, geometric constraints, and affine curvature flow, but their bibliography is absent. Source contact addresses are \texttt{rlfrank@math.princeton.edu} and \texttt{Laugesen@illinois.edu}; they are provenance, not an instruction to contact anyone.
\item Supplied research summary, ``A correction-transfer triangle for Polya, Riesz means, and normalized eigenvalue means,'' attempt 1. No full reference for the unit-square two-term asymptotic is supplied. The original literature field is empty despite references embedded in the statement and remarks.
\end{itemize}
\end{document}
